\documentclass[11pt,a4paper]{article}
\usepackage[utf8]{inputenc}
\usepackage[T1]{fontenc}
\usepackage{amsmath,amssymb,amsthm}
\usepackage{mathtools}
\usepackage{physics}
\usepackage[margin=2.5cm]{geometry}
\usepackage{hyperref}
\usepackage{cleveref}
\usepackage[shortlabels]{enumitem}
\usepackage{xcolor}
\usepackage{tcolorbox}
\usepackage{booktabs}
\usepackage{fancyhdr}
\usepackage{longtable}

% Custom commands (matching NT-4a/NT-4b/NT-4c conventions)
\newcommand{\Id}{\mathbb{1}}
\newcommand{\Ricsq}{R_{\mu\nu}R^{\mu\nu}}
\newcommand{\Rsq}{R^2}
\newcommand{\Weylsq}{C_{\mu\nu\rho\sigma}C^{\mu\nu\rho\sigma}}
\renewcommand{\dd}{\mathrm{d}}
\newcommand{\ie}{\textit{i.e.}}
\newcommand{\eg}{\textit{e.g.}}
\newcommand{\erfi}{\operatorname{erfi}}
\newcommand{\Fhat}{\hat{F}}
\newcommand{\PiTT}{\Pi_{\mathrm{TT}}}
\newcommand{\Pis}{\Pi_{\mathrm{s}}}

% Theorem environments
\theoremstyle{plain}
\newtheorem{theorem}{Theorem}[section]
\newtheorem{proposition}[theorem]{Proposition}
\newtheorem{lemma}[theorem]{Lemma}
\newtheorem{corollary}[theorem]{Corollary}
\theoremstyle{definition}
\newtheorem{definition}[theorem]{Definition}
\newtheorem{remark}[theorem]{Remark}

% Verification annotations
\newcommand{\verified}[1]{%
  \par\smallskip\noindent
  \colorbox{green!10}{\parbox{\dimexpr\linewidth-2\fboxsep}{%
    \textcolor{green!50!black}{\textbf{[Verified]}} #1}}%
  \par\smallskip}

\newcommand{\signcorrection}[1]{%
  \par\smallskip\noindent
  \colorbox{red!10}{\parbox{\dimexpr\linewidth-2\fboxsep}{%
    \textcolor{red!60!black}{\textbf{[Critical issue]}} #1}}%
  \par\smallskip}

\newcommand{\gap}[1]{%
  \par\smallskip\noindent
  \colorbox{yellow!15}{\parbox{\dimexpr\linewidth-2\fboxsep}{%
    \textcolor{orange!70!black}{\textbf{[Gap]}} #1}}%
  \par\smallskip}

% Header
\pagestyle{fancy}
\fancyhf{}
\fancyhead[L]{\small SCT Theory --- Literature Review}
\fancyhead[R]{\small INF-1: Spectral Inflation}
\fancyfoot[C]{\thepage}

\hypersetup{
  colorlinks=true,
  linkcolor=blue!60!black,
  citecolor=green!50!black,
  urlcolor=blue!50!black,
  pdfauthor={David Alfyorov},
  pdftitle={INF-1: Literature Review for Spectral Inflation},
  pdfsubject={Spectral Causal Theory --- Inflation from the R-squared sector}
}

\title{INF-1: Literature Review\\[6pt]
\large Inflationary Observables from the Spectral Action's\\
$R^2$ Sector: Conventions, Benchmarks, and Gap Analysis}
\author{David Alfyorov}
\date{March 2026}

\begin{document}
\maketitle

\begin{center}
\small\textit{Credits: Aliaksandr Samatyia contributed research-assistance and workflow support.}
\end{center}

\begin{abstract}
This document collects the conventions, formulas, benchmarks, and
references needed to derive inflationary observables ($n_s$, $r$,
$f_{\mathrm{NL}}$, $T_{\mathrm{reh}}$) from the SCT spectral action's
$R^2$ sector. We review standard Starobinsky inflation, its realization
in noncommutative geometry (NCG) via the spectral action, and the
nonlocal corrections studied by Koshelev--Kumar--Starobinsky (KKS).
We identify the connection between SCT's verified form factors
$F_1$, $F_2$ and the KKS framework, analyze the critical scalaron mass
problem, and catalog the observational constraints from Planck~2018,
BICEP/Keck~2021, and upcoming experiments.

The document serves as the literature review for the INF-1 derivation.
\end{abstract}

\tableofcontents

%======================================================================
\section{Introduction and Scope}\label{sec:intro}
%======================================================================

The spectral action $S_{\mathrm{spec}} = \mathrm{Tr}(f(D^2/\Lambda^2))$
naturally produces an $R^2$ term in the gravitational sector, with the
coefficient determined by the spectral data. In SCT, the one-loop
effective action takes the nonlocal form
\begin{equation}\label{eq:sct-action}
  S = \frac{1}{2\kappa^2}\int \dd^4 x\,\sqrt{g}\,R
  + \frac{1}{16\pi^2}\int \dd^4 x\,\sqrt{g}\,
  \Bigl[\alpha_C\,C_{\mu\nu\rho\sigma}\,F_1(\Box/\Lambda^2)\,
  C^{\mu\nu\rho\sigma}
  + \alpha_R(\xi)\,R\,F_2(\Box/\Lambda^2,\xi)\,R\Bigr],
\end{equation}
where $\alpha_C = 13/120$ and $\alpha_R(\xi) = 2(\xi - 1/6)^2$ are the
SM-summed heat kernel coefficients (NT-1b Phase~3), $F_1$ and $F_2$ are
entire form factors (NT-2), $\Lambda$ is the spectral cutoff, and
$\xi$ is the Higgs non-minimal coupling.

The $R^2$ sector drives Starobinsky-type inflation. The goal of INF-1 is
to compute the inflationary observables and compare with CMB data.

\subsection{Dependencies}
\begin{itemize}[nosep]
  \item \textbf{NT-1b (Phase~3):} $\alpha_C = 13/120$,
    $\alpha_R(\xi) = 2(\xi - 1/6)^2$. SM counting:
    $N_s = 4$, $N_D = 22.5$, $N_v = 12$.
  \item \textbf{NT-2:} $\varphi(z)$, $F_1(z)$, $F_2(z,\xi)$ entire.
    $F_2(0,\xi=0) \approx 3.518 \times 10^{-4}$.
  \item \textbf{NT-4a:} Linearized propagator:
    $\Pis(z,\xi) = 1 + 6(\xi-1/6)^2\,z\,\Fhat_2(z,\xi)$.
    Scalar mass: $m_0(\xi=0) = \sqrt{6}\,\Lambda \approx 2.449\,\Lambda$.
  \item \textbf{NT-4b:} Full nonlinear field equations in Weyl basis.
  \item \textbf{NT-4c:} FLRW reduction. Modified Friedmann equations.
    $c_T = c$. De~Sitter exact solution. Weyl sector drops out on FLRW.
\end{itemize}

%======================================================================
\section{Standard Starobinsky Inflation}\label{sec:starobinsky}
%======================================================================

\subsection{The Local $R + R^2$ Model}

The Starobinsky model~\cite{Star80} is defined by the action
\begin{equation}\label{eq:starobinsky-action}
  S = \frac{1}{2}\int \dd^4 x\,\sqrt{g}\,
  \Bigl(M_P^2\,R + f_0\,R^2\Bigr),
  \qquad f_0 = \frac{M_P^2}{6M^2},
\end{equation}
where $M$ is the scalaron mass and $M_P = 1/\sqrt{8\pi G}$ is the
reduced Planck mass. The trace equation yields the eigenvalue problem
\begin{equation}\label{eq:eigenvalue}
  \Box R = M^2 R,
\end{equation}
whose solution describes a quasi-de~Sitter inflationary phase with
$M^2 \ll R \ll M_P^2$.

\subsection{Einstein Frame and Scalaron Potential}

The conformal transformation
$\tilde{g}_{\mu\nu} = (1 + 2f_0 R/M_P^2)\,g_{\mu\nu}$ maps the
Jordan-frame action~\eqref{eq:starobinsky-action} to the Einstein frame
with a canonical scalar field (scalaron)
\begin{equation}\label{eq:scalaron}
  \varphi = \sqrt{\frac{3}{2}}\,M_P\,
  \ln\!\Bigl(1 + \frac{2f_0 R}{M_P^2}\Bigr)
  = \sqrt{\frac{3}{2}}\,M_P\,
  \ln\!\Bigl(1 + \frac{R}{3M^2}\Bigr).
\end{equation}
The scalaron potential is~\cite{Star80,Mukhanov81}
\begin{equation}\label{eq:V-starobinsky}
  \boxed{
  V(\varphi) = \frac{3M^2M_P^2}{4}
  \Bigl(1 - e^{-\sqrt{2/3}\,\varphi/M_P}\Bigr)^2.
  }
\end{equation}

\subsection{Slow-Roll Parameters and Predictions}

The slow-roll parameters evaluated at $N$ $e$-folds before the end of
inflation are~\cite{Star80,Mukhanov81,DeFelice10}
\begin{align}
  \epsilon &= \frac{M_P^2}{2}\left(\frac{V'}{V}\right)^2
  \approx \frac{3}{4N^2}, \label{eq:epsilon}\\
  \eta &= M_P^2\,\frac{V''}{V}
  \approx -\frac{1}{N}. \label{eq:eta}
\end{align}
The resulting inflationary observables are:
\begin{align}
  n_s &= 1 - 6\epsilon + 2\eta \approx 1 - \frac{2}{N},
  \label{eq:ns-star}\\
  r &= 16\epsilon \approx \frac{12}{N^2},
  \label{eq:r-star}\\
  n_t &= -2\epsilon \approx -\frac{3}{2N^2}.
  \label{eq:nt-star}
\end{align}

\begin{center}
\renewcommand{\arraystretch}{1.3}
\begin{tabular}{lcc}
\toprule
Observable & $N = 50$ & $N = 60$ \\
\midrule
$n_s$ & 0.960 & 0.967 \\
$r$ & $4.8 \times 10^{-3}$ & $3.3 \times 10^{-3}$ \\
$n_t$ & $-6.0 \times 10^{-4}$ & $-4.2 \times 10^{-4}$ \\
\bottomrule
\end{tabular}
\end{center}

The single-field consistency relation $r = -8n_t$ is satisfied.

\subsection{Amplitude Normalization}

The scalar power spectrum amplitude fixes $M$ via~\cite{Mukhanov81}
\begin{equation}\label{eq:As}
  \mathcal{P}_{\mathcal{R}} = \frac{N^2 M^2}{24\pi^2 M_P^2}
  = A_s \approx 2.1 \times 10^{-9}
  \qquad\Longrightarrow\qquad
  M \approx 1.3 \times 10^{-5}\,M_P \approx 3.1 \times 10^{13}\;\text{GeV}.
\end{equation}

\verified{Standard Starobinsky predictions verified numerically:
$n_s(55) = 0.9636$, $r(55) = 3.97 \times 10^{-3}$,
$M/M_P = 1.28 \times 10^{-5}$.}

%======================================================================
\section{Starobinsky Inflation in NCG / Spectral Action}
\label{sec:ncg}
%======================================================================

\subsection{The Spectral Action's $R^2$ Term}

In the noncommutative geometry approach of Chamseddine and
Connes~\cite{CC97,CC06}, the spectral action
$S = \mathrm{Tr}(f(D^2/\Lambda^2))$ produces, in its asymptotic
expansion, the gravitational terms
\begin{equation}\label{eq:spectral-expansion}
  S_{\mathrm{spec}} \sim \frac{f_2\Lambda^2}{4\pi^2}\int \dd^4 x\,\sqrt{g}\,R
  + \frac{f_4}{320\pi^2}\int \dd^4 x\,\sqrt{g}\,
  \Bigl(-3\,\Weylsq + \frac{11}{6}\,R^*R^*
  + 10\,\Rsq\Bigr) + \cdots
\end{equation}
where $f_n = \int_0^\infty f(u)\,u^{n/2-1}\dd u$ are moments of the
cutoff function and $R^*R^*$ is the Gauss--Bonnet density. The $R^2$
coefficient is $c_2 = f_4/(16\pi^2)$ in the minimal spectral triple.

In the full SM spectral triple, the matter sector contributes through
the heat kernel one-loop effective action~\cite{CC12,CCS13}, producing
exactly the structure computed in SCT's NT-1b:
\begin{equation}\label{eq:c2-sct}
  c_2^{\mathrm{total}} = \frac{\alpha_R(\xi)}{16\pi^2}
  = \frac{2(\xi - 1/6)^2}{16\pi^2}.
\end{equation}

\subsection{Conformal Coupling and the Inflation Problem}

\signcorrection{At conformal coupling $\xi = 1/6$:
$\alpha_R(1/6) = 0$ and $c_2 = 0$.
There is \textbf{no $R^2$ term} and therefore \textbf{no Starobinsky
inflation} from the one-loop spectral action alone. This is a
fundamental constraint: inflation in SCT requires $\xi \neq 1/6$.}

\subsection{NCG Inflation Literature}

\paragraph{Chamseddine--Connes (2006, 2012)~\cite{CC06,CC12}.}
The spectral action naturally gives rise to an $R^2$ term that can
drive Starobinsky inflation. The coefficients are determined by the
spectral data.

\paragraph{Nelson--Sakellariadou (2009)~\cite{NS09}.}
Demonstrated that inflation arises naturally in the NCG approach
through the non-minimal coupling between geometry and the Higgs field.
The inflationary phase occurs without introducing additional
fields---it is a consequence of the $\xi$-dependent coupling.

\paragraph{Buck--Fairbairn--Sakellariadou (2010)~\cite{BFS10}.}
Analyzed slow-roll inflation with conformally coupled inflaton in the
spectral action context. Found that Higgs inflation in NCG is
incompatible with the measured top quark mass due to CMB constraints
on the self-coupling running. Analyzed the role of additional
conformally coupled massless scalars arising in NCG.

\paragraph{Chamseddine--Connes--van~Suijlekom (2013)~\cite{CCS13}.}
Connected the spectral action directly to Starobinsky inflation,
emphasizing that the $R^2$ coefficient is determined by the geometry
of the spectral triple.

\paragraph{Devastato--Lizzi--Martinetti (2014)~\cite{DLM14}.}
Studied the ``grand symmetry'' of the spectral action and its
implications for the Higgs mass and inflation in the NCG framework.

%======================================================================
\section{Scalaron Mass from Spectral Data: Critical Analysis}
\label{sec:scalaron-mass}
%======================================================================

\subsection{The Mass Formula}

In the local limit ($F_2 \to c_2$), the spectral action's $R^2$ sector
gives a Starobinsky model with scalaron mass
\begin{equation}\label{eq:M-from-c2}
  M^2 = \frac{M_P^2}{12\,c_2^{\mathrm{total}}}
  = \frac{M_P^2 \cdot 16\pi^2}{12 \cdot 2(\xi - 1/6)^2}
  = \frac{2\pi^2 M_P^2}{3(\xi - 1/6)^2}.
\end{equation}

Equivalently, using the NT-4a scalar propagator pole:
\begin{equation}\label{eq:m0-from-nt4a}
  m_0^2(\xi) = \frac{\Lambda^2}{6(\xi - 1/6)^2},
  \qquad m_0(\xi = 0) = \sqrt{6}\,\Lambda \approx 2.449\,\Lambda.
\end{equation}

\subsection{The Scalaron Mass Problem}

\signcorrection{
For $\xi = 0$ (minimal coupling):
\begin{equation}\label{eq:M-xi0}
  M = \sqrt{\frac{2\pi^2}{3 \cdot (1/6)^2}}\,M_P
  = \sqrt{24\pi^2}\,M_P \approx 15.4\,M_P.
\end{equation}
\textbf{This is far too heavy for inflation.} The observed amplitude
$A_s \approx 2.1 \times 10^{-9}$ requires
$M \approx 1.3 \times 10^{-5}\,M_P$.

The scalaron mass from the pure SM spectral action at $\xi = 0$
exceeds the required value by a factor of $\sim 1.2 \times 10^6$.
This is the central challenge for spectral inflation in SCT.}

\subsection{Resolution Paths}

There are several possible resolutions:

\begin{enumerate}[(a)]
  \item \textbf{Nonlocal form factors.}
    In the full nonlocal theory, the scalaron mass is modified by
    $F_2(\Box/\Lambda^2)$. The KKS framework~\cite{KKS17,KKS20,KKS22}
    shows that the scalar power spectrum is preserved while the tensor
    spectrum is modified. However, the \emph{background} inflationary
    solution $\Box R = M^2 R$ is identical to the local $R^2$ case.
    The effective scalaron mass at the background level remains
    $M^2 = M_P^2/(12c_2)$, so the mass problem is not resolved by
    nonlocality alone.

  \item \textbf{Cutoff scale $\Lambda \ll M_P$.}
    If $\Lambda \sim 1.3 \times 10^{13}$~GeV $(\approx 5 \times 10^{-6}\,M_P)$
    with $\xi = 0$, then
    $m_0 = \sqrt{6}\,\Lambda \approx 3.1 \times 10^{13}$~GeV
    $\approx 1.3 \times 10^{-5}\,M_P$, matching the required scalaron
    mass. This is a viable option but requires $\Lambda$ to be
    sub-Planckian by five orders of magnitude.
    \emph{Note:} This identification of the NT-4a pole mass $m_0$ with
    the Starobinsky scalaron mass $M$ is valid in the full spectral
    action (where $M_P$ itself depends on $\Lambda$ through
    $M_P^2 = f_2 \Lambda^2/(2\pi^2)$); in the one-loop EFT approach
    with $M_P$ as an external input, $M$ is determined by $c_2$ and
    $M_P$ independently of $\Lambda$.

  \item \textbf{Large non-minimal coupling $\xi$.}
    With $\Lambda = M_P$ and arbitrary $\xi$:
    \[
      M \approx 1.3 \times 10^{-5}\,M_P
      \quad\Longrightarrow\quad
      |\xi - 1/6| \approx 3.2 \times 10^{4}.
    \]
    This enormous value of $\xi$ is unnatural and violates perturbative
    control of the non-minimal coupling.

  \item \textbf{Additional scalar fields.}
    The NCG spectral triple may contain additional scalars (e.g., the
    $\sigma$-field from the seesaw sector, a dilaton from scale
    invariance~\cite{CC05}). These would increase $\alpha_R$ and lower
    the scalaron mass. This is the approach explored in the NCG
    inflation literature~\cite{NS09,BFS10}.

  \item \textbf{Two-loop or higher-order corrections.}
    The one-loop result gives only the leading contribution to $c_2$.
    Higher-loop corrections could modify $\alpha_R$ substantially if
    the theory is near a fixed point.
\end{enumerate}

\gap{G1: Determine the viable parameter space $(\Lambda, \xi)$ for
which the SCT spectral action produces the correct $A_s$.
This requires solving $M(\Lambda, \xi) = 1.3 \times 10^{-5}\,M_P$
subject to the PPN-1 constraint $\Lambda \geq 2.38 \times 10^{-3}$~eV.}

%======================================================================
\section{Nonlocal Corrections: The KKS Framework}\label{sec:kks}
%======================================================================

\subsection{The Generalized Nonlocal $R^2$ Action}

Koshelev, Kumar, and Starobinsky~\cite{KKS16,KKS17,KKS20,KKS22,KKS22b}
developed the most comprehensive framework for inflation in nonlocal
gravity with quadratic curvature terms. Their quadratic-in-curvature
action reads
\begin{equation}\label{eq:kks-action}
  S = \int \dd^4 x\,\sqrt{-g}\,
  \Bigl(\frac{M_P^2}{2}\,R + R\,\mathcal{F}_R(\Box_s)\,R
  + W_{\mu\nu\rho\sigma}\,\mathcal{F}_C(\Box_s)\,
  W^{\mu\nu\rho\sigma}\Bigr),
\end{equation}
where $\Box_s = \Box/\mathcal{M}_s^2$ and $\mathcal{M}_s$ is the
nonlocality scale. The form factors are constrained to be
\begin{align}
  \mathcal{F}_R(\Box_s) &= f_0 M^2\,
  \frac{1 - (1 - \Box/M^2)\,e^{\gamma_S(\Box_s)}}{\Box},
  \label{eq:FR}\\
  \mathcal{F}_C(\Box_s) &= \frac{M_P^2}{2}\,
  \frac{e^{\gamma_T(\Box_s)} - 1}{\Box},
  \label{eq:FC}
\end{align}
where $\gamma_S$, $\gamma_T$ are entire functions (ensuring ghost-freedom).

\subsection{Key Results for Inflation}

The KKS framework establishes~\cite{KKS16,KKS17,KKS22}:

\begin{enumerate}[(1)]
  \item \textbf{Background solution.}
    $\Box R = M^2 R$ (the Starobinsky eigenvalue equation) is an
    \emph{exact} background solution for any choice of $\gamma_S$,
    $\gamma_T$ satisfying $\mathcal{F}_R(M^2/\mathcal{M}_s^2) = f_0$
    and $\mathcal{F}_R'(M^2/\mathcal{M}_s^2) = 0$.

  \item \textbf{Scalar power spectrum.}
    $\mathcal{P}_{\mathcal{R}}$ and $n_s$ are \textbf{unchanged}
    from local $R^2$ inflation:
    \begin{equation}\label{eq:ns-kks}
      n_s \approx 1 - \frac{2}{N}.
    \end{equation}
    Nonlocal corrections to $n_s$ are $O(M^4/\mathcal{M}_s^4)$,
    negligible for $M \ll \mathcal{M}_s$.

  \item \textbf{Tensor power spectrum.}
    Modified by the Weyl form factor:
    \begin{equation}\label{eq:Pt-kks}
      \mathcal{P}_T = \frac{1}{12\pi^2 f_0}
      (1 - 3\epsilon)\,e^{-2\gamma_T(-\bar{R}_{\mathrm{dS}}/(6\mathcal{M}_s^2))}.
    \end{equation}

  \item \textbf{Tensor-to-scalar ratio.}
    \begin{equation}\label{eq:r-kks}
      r = \frac{12}{N^2}\,e^{-2\gamma_T(-\bar{R}_{\mathrm{dS}}/(6\mathcal{M}_s^2))}.
    \end{equation}
    The exponential factor can suppress or enhance $r$ relative to the
    local Starobinsky value. Any $r < 0.036$ is achievable.

  \item \textbf{Modified consistency relation.}
    $r \neq -8n_t$ in general. The tensor tilt can be positive
    ($n_t > 0$, blue tilt) or negative.

  \item \textbf{Sound speed.}
    $c_s^2 = c_T^2 = 1$ (both scalar and tensor). Compatible with
    GW170817.

  \item \textbf{Degrees of freedom.}
    Around Minkowski: scalaron (massive spin-0) + graviton (massless
    spin-2). No ghosts for entire $\gamma_{S,T}$.
\end{enumerate}

\subsection{Connection to SCT Form Factors}

The SCT action~\eqref{eq:sct-action} maps to the KKS
framework~\eqref{eq:kks-action} via the identification:
\begin{equation}\label{eq:sct-kks-map}
  \boxed{
  \mathcal{F}_R(\Box_s) = \frac{\alpha_R(\xi)}{16\pi^2}\,
  F_2(\Box/\Lambda^2,\xi),
  \qquad
  \mathcal{F}_C(\Box_s) = \frac{\alpha_C}{16\pi^2}\cdot 2\,
  F_1(\Box/\Lambda^2),
  \qquad
  \mathcal{M}_s = \Lambda.
  }
\end{equation}

The SCT form factors $F_1$, $F_2$ are entire (NT-2), which is
precisely the ghost-freedom condition required in the KKS framework.
The difference is that SCT \textbf{derives} the form factors from the
one-loop heat kernel, while KKS treats $\gamma_S$, $\gamma_T$ as
phenomenological inputs.

\gap{G2: Verify explicitly that the SCT form factors
$\mathcal{F}_R$, $\mathcal{F}_C$ satisfy the KKS conditions
$\mathcal{F}_R(M^2/\Lambda^2) = f_0$ and
$\mathcal{F}_R'(M^2/\Lambda^2) = 0$.
If they do not, determine what modification is needed.}

%======================================================================
\section{Non-Gaussianity}\label{sec:fNL}
%======================================================================

\subsection{Single-Field Slow-Roll Prediction}

In standard single-field slow-roll inflation, the non-Gaussianity
parameter satisfies~\cite{Maldacena03,Acquaviva03}
\begin{equation}\label{eq:fNL-slowroll}
  f_{\mathrm{NL}}^{\mathrm{local}} = \frac{5}{12}(n_s - 1)
  \approx -\frac{5}{12}\cdot\frac{2}{N} \approx -0.015
  \quad\text{for } N = 55.
\end{equation}
This is the Maldacena consistency relation. Planck~2018 measures
$f_{\mathrm{NL}}^{\mathrm{local}} = -0.9 \pm 5.1$~\cite{Planck18IX}.

\subsection{Nonlocal Corrections to $f_{\mathrm{NL}}$}

Koshelev, Kumar, and Starobinsky~\cite{KKS20,KKS22b} computed
non-Gaussianities in the nonlocal $R^2$ framework:
\begin{itemize}[nosep]
  \item $|f_{\mathrm{NL}}| \sim O(1\text{--}10)$ is achievable in
    equilateral, orthogonal, and squeezed configurations.
  \item The sign of $f_{\mathrm{NL}}$ depends on the choice of form
    factors.
  \item Running of $f_{\mathrm{NL}}$:
    $|\dd \ln f_{\mathrm{NL}}/\dd \ln k| \lesssim 1$, which can be an
    order of magnitude larger than in local EFT models.
  \item Cubic-in-curvature nonlocal terms contribute to
    the bispectrum but not to two-point functions.
\end{itemize}

In SCT, the form factors are fixed by the spectral data, so
$f_{\mathrm{NL}}$ is in principle \textbf{predicted} (not a free
parameter). The cubic terms in the BV alpha-insertion
$\Theta^{(R)}_{\mu\nu}$ provide the SCT analog of the KKS cubic
curvature operators.

\gap{G3: Compute $f_{\mathrm{NL}}$ for the specific SCT form factors.
This requires the three-point correlator of curvature perturbations,
which involves the cubic terms in the spectral action expansion.}

%======================================================================
\section{Reheating}\label{sec:reheating}
%======================================================================

\subsection{Scalaron Decay}

After inflation, the scalaron oscillates around the minimum of the
potential~\eqref{eq:V-starobinsky} and decays into SM particles. The
dominant decay channel is into pairs of conformally coupled
fields~\cite{Star83,Vilenkin85}:
\begin{equation}\label{eq:decay-rate}
  \Gamma_\varphi \sim \frac{M^3}{M_P^2}.
\end{equation}
The reheating temperature is estimated as
\begin{equation}\label{eq:Treh}
  T_{\mathrm{reh}} \sim \Bigl(\frac{90}{\pi^2 g_*}\Bigr)^{1/4}
  \sqrt{\Gamma_\varphi\,M_P}
  \sim \left(\frac{M^3}{M_P}\right)^{1/2}
  \approx 6 \times 10^{10}\;\text{GeV}
\end{equation}
for $M \approx 3.1 \times 10^{13}$~GeV and
$g_* = 106.75$ (SM degrees of freedom).

\subsection{BBN Constraint}

The BBN constraint requires~\cite{Kawasaki00}
\begin{equation}\label{eq:bbn}
  T_{\mathrm{reh}} > 5\;\text{MeV}.
\end{equation}
Starobinsky inflation satisfies this by many orders of magnitude.

\subsection{Reheating $e$-Folds and $n_s$}

The number of $e$-folds between Hubble exit and the end of inflation
depends on the reheating history:
\begin{equation}\label{eq:N-reh}
  N = 55 - \frac{1}{4}\ln\!\left(\frac{10^{16}\;\text{GeV}}{V_{\mathrm{end}}^{1/4}}\right)
  + \frac{1}{4}\ln\!\left(\frac{T_{\mathrm{reh}}}{10^9\;\text{GeV}}\right).
\end{equation}
For Starobinsky inflation with efficient reheating,
$N \approx 50\text{--}60$.

\gap{G4: In the nonlocal theory, the scalaron coupling to matter
is modified by the form factors. Compute the modified decay rate
$\Gamma_\varphi$ and verify $T_{\mathrm{reh}} > 5$~MeV.}

%======================================================================
\section{CMB Observational Constraints}\label{sec:cmb}
%======================================================================

\subsection{Current Data}

\begin{center}
\renewcommand{\arraystretch}{1.3}
\begin{tabular}{lcc}
\toprule
Observable & Value & Reference \\
\midrule
$n_s$ (TT+TE+EE+lowE+lensing)
  & $0.9649 \pm 0.0042$ (68\%~CL) & Planck 2018~\cite{Planck18X} \\
$r$ (95\%~CL)
  & $< 0.036$ & BICEP/Keck 2021~\cite{BK21} \\
$f_{\mathrm{NL}}^{\mathrm{local}}$
  & $-0.9 \pm 5.1$ & Planck 2018~\cite{Planck18IX} \\
$f_{\mathrm{NL}}^{\mathrm{equil}}$
  & $-26 \pm 47$ & Planck 2018~\cite{Planck18IX} \\
$\dd n_s/\dd \ln k$
  & $-0.0045 \pm 0.0067$ & Planck 2018~\cite{Planck18X} \\
\bottomrule
\end{tabular}
\end{center}

\subsection{Future Experiments}

\begin{center}
\renewcommand{\arraystretch}{1.3}
\begin{tabular}{lccc}
\toprule
Experiment & $\sigma(r)$ & $\sigma(n_s)$ & Timeline \\
\midrule
CMB-S4 & $\sim 10^{-3}$ & $\sim 0.002$ & 2030s \\
LiteBIRD & $\sim 10^{-3}$ & $\sim 0.002$ & 2030s \\
Simons Observatory & $\sim 0.003$ & $\sim 0.003$ & 2020s \\
\bottomrule
\end{tabular}
\end{center}

\textbf{Detection prospect:} Standard Starobinsky predicts
$r \approx 3.3 \times 10^{-3}$ for $N = 60$, which is within the
$\sim 3\sigma$ reach of LiteBIRD and CMB-S4. If nonlocal corrections
suppress $r$ (as allowed in the KKS framework), detection becomes
more challenging.

%======================================================================
\section{De Sitter Conjecture and Swampland}\label{sec:swampland}
%======================================================================

The refined de~Sitter conjecture~\cite{OPSV18} states that in any
consistent quantum gravity theory:
\begin{equation}\label{eq:dS-conjecture}
  |\nabla V|/V \geq c \sim O(1)
  \qquad\text{or}\qquad
  \min(\nabla_i\nabla_j V) \leq -c'\,V.
\end{equation}
The Starobinsky potential~\eqref{eq:V-starobinsky} satisfies the
\emph{second} condition (the mass matrix has a negative eigenvalue at
the inflection point) but not the first (the plateau violates
$|\nabla V|/V \geq c$).

\gap{G5: Evaluate the Swampland criteria for the SCT
potential explicitly. Determine whether the nonlocal corrections
modify the Einstein-frame potential shape enough to affect the
de~Sitter conjecture status.}

%======================================================================
\section{NT-4c Results Relevant to INF-1}\label{sec:nt4c}
%======================================================================

From the verified NT-4c FLRW reduction~\cite{NT4c}:

\begin{enumerate}[(1)]
  \item \textbf{Weyl sector drops out on FLRW.}
    $C_{\mu\nu\rho\sigma} = 0$ eliminates $F_1$, $B_{\mu\nu}$,
    $\Theta^{(C)}_{\mu\nu}$ from the \emph{background} equations.
    Only the $R^2$ sector (governed by $\alpha_R(\xi)$) contributes.

  \item \textbf{Modified Friedmann equations.}
    \begin{equation}\label{eq:friedmann-recap}
      H^2 = \frac{\kappa^2}{3}\,\rho
      - \frac{\kappa^2}{3}\,\beta_R\,
      [H_{00} + \Theta^{(R)}_{00}],
      \qquad
      \beta_R = \frac{\alpha_R(\xi)}{16\pi^2}.
    \end{equation}

  \item \textbf{GW speed: $c_T = c$.}
    The $R^2$ sector couples only to scalar perturbations (SVT
    decomposition). Tensor mode speed is unmodified. This is consistent
    with the KKS result $c_T^2 = 1$ and the GW170817 constraint.

  \item \textbf{De~Sitter is exact solution.}
    $H_{\mu\nu}|_{\mathrm{dS}} = 0$,
    $\Theta^{(R)}_{\mu\nu}|_{\mathrm{dS}} = 0$.

  \item \textbf{Alpha-insertion EOS: $w_\Theta = +1$ (stiff).}
    The nonlocal $\Theta^{(R)}$ acts as kinetic energy, not vacuum
    energy. This is important for the inflationary dynamics: the
    alpha-insertion contributes during the oscillatory reheating phase,
    not during slow roll.

  \item \textbf{Conformal coupling ($\xi = 1/6$): no inflation.}
    All spectral corrections vanish identically.

  \item \textbf{Scalaron equation from trace.}
    The trace of the modified field equations, under the ansatz
    $\Box R = M^2 R$, yields the Starobinsky scalaron equation.
    This connects the NT-4c modified Friedmann equations to the
    KKS framework.
\end{enumerate}

%======================================================================
\section{Conventions and Notation Summary}\label{sec:conventions}
%======================================================================

\begin{center}
\renewcommand{\arraystretch}{1.3}
\begin{longtable}{ll}
\toprule
Symbol & Definition \\
\midrule
$g_{\mu\nu}$ & Lorentzian metric, signature $(-,+,+,+)$ \\
$\kappa^2 = 8\pi G = M_P^{-2}$ & Gravitational coupling \\
$M_P = 2.435 \times 10^{18}$~GeV & Reduced Planck mass \\
$\Lambda$ & Spectral cutoff scale \\
$\xi$ & Higgs non-minimal coupling to gravity \\
$\alpha_C = 13/120$ & SM Weyl$^2$ coefficient ($\xi$-independent) \\
$\alpha_R(\xi) = 2(\xi - 1/6)^2$ & SM $R^2$ coefficient \\
$c_2 = \alpha_R/(16\pi^2)$ & Local $R^2$ coupling \\
$M$ & Scalaron mass: $M^2 = M_P^2/(12c_2)$ \\
$f_0 = M_P^2/(6M^2)$ & Starobinsky coefficient \\
$m_0(\xi)$ & NT-4a scalar pole: $m_0 = \Lambda/\sqrt{6(\xi-1/6)^2}$ \\
$m_2$ & NT-4a tensor pole: $m_2 = \Lambda\sqrt{60/13}$ \\
$\varphi(z) = \int_0^1 e^{-\xi(1-\xi)z}\,\dd\xi$ & Master function \\
$F_1(z)$, $F_2(z,\xi)$ & Entire form factors (NT-2) \\
$\Fhat_i(z) = F_i(z)/F_i(0)$ & Normalized form factors \\
$N$ & Number of $e$-folds before end of inflation \\
$n_s$ & Scalar spectral index \\
$r$ & Tensor-to-scalar ratio \\
$n_t$ & Tensor spectral index \\
$f_{\mathrm{NL}}$ & Non-Gaussianity parameter \\
$T_{\mathrm{reh}}$ & Reheating temperature \\
$A_s \approx 2.1 \times 10^{-9}$ & Scalar power spectrum amplitude \\
\bottomrule
\end{longtable}
\end{center}

%======================================================================
\section{Gap Analysis and Derivation Plan}\label{sec:gaps}
%======================================================================

\begin{center}
\renewcommand{\arraystretch}{1.4}
\begin{tabular}{clp{7cm}c}
\toprule
\# & Gap & Description & Severity \\
\midrule
G1 & Scalaron mass & Determine viable $(\Lambda, \xi)$ parameter
  space for $A_s$ matching & \textbf{CRITICAL} \\
G2 & KKS conditions & Verify SCT form factors satisfy
  $\mathcal{F}_R(M^2/\Lambda^2) = f_0$,
  $\mathcal{F}_R'(M^2/\Lambda^2) = 0$ & HIGH \\
G3 & $f_{\mathrm{NL}}$ & Compute non-Gaussianity for SCT form
  factors (requires bispectrum) & MEDIUM \\
G4 & Reheating & Compute modified decay rate from nonlocal
  couplings & MEDIUM \\
G5 & Swampland & Evaluate de~Sitter conjecture for SCT
  Einstein-frame potential & LOW \\
G6 & Perturbation theory & Derive full second-order perturbation
  action around FLRW in SCT (tensor + scalar) & HIGH \\
G7 & Weyl sector on tensors & The Weyl sector drops out of
  \emph{background} equations but contributes to tensor
  \emph{perturbations}. Quantify this effect. & HIGH \\
\bottomrule
\end{tabular}
\end{center}

\subsection{Derivation Strategy for INF-1}

\begin{enumerate}[(1)]
  \item \textbf{Local limit first.}
    Compute all observables in the local Starobinsky limit ($F_2 \to c_2$)
    as a function of $(\Lambda, \xi)$. Determine the viable parameter
    space from $A_s$ matching and PPN-1 constraints.

  \item \textbf{Nonlocal tensor corrections.}
    Using the KKS framework with SCT form factors, compute the modified
    tensor power spectrum and the resulting $r$, $n_t$.

  \item \textbf{Scalar perturbation check.}
    Verify that scalar perturbations are unchanged from the local case
    (as predicted by KKS).

  \item \textbf{Reheating.}
    Compute $\Gamma_\varphi$ and $T_{\mathrm{reh}}$ in the local limit.
    Estimate nonlocal corrections.

  \item \textbf{Non-Gaussianity (if feasible).}
    Compute $f_{\mathrm{NL}}$ from the SCT cubic terms, or bound it.

  \item \textbf{Comparison.}
    Tabulate $(n_s, r, n_t, f_{\mathrm{NL}}, T_{\mathrm{reh}})$
    for SCT vs.\ standard Starobinsky vs.\ KKS nonlocal Starobinsky
    vs.\ Planck/BICEP data.
\end{enumerate}

%======================================================================
\begin{thebibliography}{99}
%======================================================================

\bibitem{Star80}
A.~A.~Starobinsky,
``A new type of isotropic cosmological models without singularity,''
\emph{Phys.\ Lett.\ B} \textbf{91} (1980) 99--102.

\bibitem{Mukhanov81}
V.~F.~Mukhanov and G.~V.~Chibisov,
``Quantum fluctuation and `nonsingular' universe,''
\emph{JETP Lett.} \textbf{33} (1981) 532--535.

\bibitem{Star83}
A.~A.~Starobinsky,
``The perturbation spectrum evolving from a nonsingular initially
de~Sitter cosmology and the microwave background anisotropy,''
\emph{Sov.\ Astron.\ Lett.} \textbf{9} (1983) 302--304.

\bibitem{Vilenkin85}
A.~Vilenkin,
``Classical and quantum cosmology of the Starobinsky inflationary model,''
\emph{Phys.\ Rev.\ D} \textbf{32} (1985) 2511.

\bibitem{CC97}
A.~H.~Chamseddine and A.~Connes,
``The spectral action principle,''
\emph{Commun.\ Math.\ Phys.} \textbf{186} (1997) 731--750,
\href{https://arxiv.org/abs/hep-th/9606001}{hep-th/9606001}.

\bibitem{CC05}
A.~H.~Chamseddine and A.~Connes,
``Scale invariance in the spectral action,''
\href{https://arxiv.org/abs/hep-th/0512169}{hep-th/0512169}.

\bibitem{CC06}
A.~H.~Chamseddine and A.~Connes,
``Inner fluctuations of the spectral action,''
\emph{J.\ Geom.\ Phys.} \textbf{57} (2006) 1--21,
\href{https://arxiv.org/abs/hep-th/0512169}{hep-th/0512169}.

\bibitem{CC12}
A.~H.~Chamseddine and A.~Connes,
``Resilience of the spectral standard model,''
\emph{JHEP} \textbf{1209} (2012) 104,
\href{https://arxiv.org/abs/1208.1030}{arXiv:1208.1030}.

\bibitem{CCS13}
A.~H.~Chamseddine, A.~Connes, and W.~D.~van~Suijlekom,
``Beyond the spectral standard model: emergence of Pati--Salam
unification,''
\emph{JHEP} \textbf{1311} (2013) 132,
\href{https://arxiv.org/abs/1304.8050}{arXiv:1304.8050}.

\bibitem{NS09}
W.~Nelson and M.~Sakellariadou,
``Natural inflation mechanism in asymptotic noncommutative geometry,''
\emph{Phys.\ Lett.\ B} \textbf{680} (2009) 263--266,
\href{https://arxiv.org/abs/0903.1520}{arXiv:0903.1520}.

\bibitem{BFS10}
M.~Buck, M.~Fairbairn, and M.~Sakellariadou,
``Inflation in models with conformally coupled scalar fields:
an application to the noncommutative spectral action,''
\emph{Phys.\ Rev.\ D} \textbf{82} (2010) 043509,
\href{https://arxiv.org/abs/1005.1188}{arXiv:1005.1188}.

\bibitem{DLM14}
A.~Devastato, F.~Lizzi, and P.~Martinetti,
``Grand symmetry, spectral action and the Higgs mass,''
\emph{JHEP} \textbf{1401} (2014) 042,
\href{https://arxiv.org/abs/1304.0415}{arXiv:1304.0415}.

\bibitem{KKS16}
A.~S.~Koshelev, L.~Modesto, L.~Rachwa{\l}, and A.~A.~Starobinsky,
``Occurrence of exact $R^2$ inflation in non-local UV-complete gravity,''
\emph{JHEP} \textbf{1611} (2016) 067,
\href{https://arxiv.org/abs/1604.03127}{arXiv:1604.03127}.

\bibitem{KKS17}
A.~S.~Koshelev, K.~S.~Kumar, and A.~A.~Starobinsky,
``$R^2$ inflation to probe non-perturbative quantum gravity,''
\emph{JHEP} \textbf{1803} (2018) 071,
\href{https://arxiv.org/abs/1711.08864}{arXiv:1711.08864}.

\bibitem{KKS20}
A.~S.~Koshelev, K.~S.~Kumar, A.~Mazumdar, and A.~A.~Starobinsky,
``Non-Gaussianities and tensor-to-scalar ratio in non-local
$R^2$-like inflation,''
\emph{JHEP} \textbf{2006} (2020) 152,
\href{https://arxiv.org/abs/2003.00629}{arXiv:2003.00629}.

\bibitem{KKS22}
A.~S.~Koshelev, K.~S.~Kumar, and A.~A.~Starobinsky,
``Generalized non-local $R^2$-like inflation,''
\href{https://arxiv.org/abs/2209.02515}{arXiv:2209.02515}.

\bibitem{KKS22b}
A.~S.~Koshelev, K.~S.~Kumar, and A.~A.~Starobinsky,
``Non-Gaussianities in generalized non-local $R^2$-like inflation,''
\href{https://arxiv.org/abs/2210.16459}{arXiv:2210.16459}.

\bibitem{KKS23}
A.~S.~Koshelev, K.~S.~Kumar, and A.~A.~Starobinsky,
``Cosmology in nonlocal gravity,''
\href{https://arxiv.org/abs/2305.18716}{arXiv:2305.18716}.

\bibitem{KST22}
A.~S.~Koshelev, A.~A.~Starobinsky, and A.~Tokareva,
``Post-inflationary GW production in generic higher (infinite)
derivative gravity,''
\href{https://arxiv.org/abs/2211.02070}{arXiv:2211.02070}.

\bibitem{BSM22}
J.~Bezerra-Sobrinho and L.~G.~Medeiros,
``Modified Starobinsky inflation by the $R\ln(\Box)R$ term,''
\href{https://arxiv.org/abs/2202.13308}{arXiv:2202.13308}.

\bibitem{Maldacena03}
J.~Maldacena,
``Non-Gaussian features of primordial fluctuations in single field
inflationary models,''
\emph{JHEP} \textbf{0305} (2003) 013,
\href{https://arxiv.org/abs/astro-ph/0210603}{astro-ph/0210603}.

\bibitem{Acquaviva03}
V.~Acquaviva, N.~Bartolo, S.~Matarrese, and A.~Riotto,
``Gauge-invariant second-order perturbations and non-Gaussianity
from inflation,''
\emph{Nucl.\ Phys.\ B} \textbf{667} (2003) 119--148.

\bibitem{DeFelice10}
A.~De~Felice and S.~Tsujikawa,
``$f(R)$ theories,''
\emph{Living Rev.\ Rel.} \textbf{13} (2010) 3,
\href{https://arxiv.org/abs/1002.4928}{arXiv:1002.4928}.

\bibitem{Kawasaki00}
M.~Kawasaki, K.~Kohri, and N.~Sugiyama,
``Cosmological constraints on late-time entropy production,''
\emph{Phys.\ Rev.\ Lett.} \textbf{82} (1999) 4168.

\bibitem{Planck18X}
Planck Collaboration,
``Planck 2018 results. X. Constraints on inflation,''
\emph{Astron.\ Astrophys.} \textbf{641} (2020) A10,
\href{https://arxiv.org/abs/1807.06211}{arXiv:1807.06211}.

\bibitem{Planck18IX}
Planck Collaboration,
``Planck 2018 results. IX. Constraints on primordial
non-Gaussianity,''
\emph{Astron.\ Astrophys.} \textbf{641} (2020) A9,
\href{https://arxiv.org/abs/1905.05697}{arXiv:1905.05697}.

\bibitem{BK21}
BICEP/Keck Collaboration,
``Improved constraints on primordial gravitational waves using
Planck, WMAP, and BICEP/Keck observations through the 2018
observing season,''
\emph{Phys.\ Rev.\ Lett.} \textbf{127} (2021) 151301,
\href{https://arxiv.org/abs/2110.00483}{arXiv:2110.00483}.

\bibitem{OPSV18}
H.~Ooguri, E.~Palti, G.~Shiu, and C.~Vafa,
``Distance and de Sitter conjectures on the Swampland,''
\emph{Phys.\ Lett.\ B} \textbf{788} (2019) 180,
\href{https://arxiv.org/abs/1810.05506}{arXiv:1810.05506}.

\bibitem{CPR09}
A.~Codello, R.~Percacci, and C.~Rahmede,
``Investigating the ultraviolet properties of gravity with a
Wilsonian renormalization group equation,''
\emph{Ann.\ Phys.} \textbf{324} (2009) 414--469,
\href{https://arxiv.org/abs/0805.2909}{arXiv:0805.2909}.

\bibitem{Ste77}
K.~S.~Stelle,
``Renormalization of higher-derivative quantum gravity,''
\emph{Phys.\ Rev.\ D} \textbf{16} (1977) 953--969.

\bibitem{NT4c}
D.~Alfyorov,
``NT-4c: FLRW reduction of SCT nonlocal field equations,''
SCT Theory internal document (2026).

\end{thebibliography}

\end{document}
